\documentclass[
	fontsize=12pt,
	paper=a4,
	numbers=noenddot,
	parskip=half,
	openany
]{scrbook}
\usepackage{tcolorbox}

\include{../../../for_latex/preamble}
% ##################################################
% Gemeinsames Setup fuer die Probeklausuren
% (Java-Listings, Hinweis-/Beispiel-Boxen, Struktogramm-Stile in TikZ)
% wird nach \include{../for_latex/preamble} per \input eingebunden
% ##################################################

% --- Java Code-Highlighting ---
\definecolor{codegray}{rgb}{0.5,0.5,0.5}
\definecolor{codepurple}{rgb}{0.58,0,0.82}
\definecolor{backcolour}{rgb}{0.95,0.95,0.92}
\definecolor{keywordcolor}{rgb}{0.8,0.47,0.13}

\lstdefinestyle{javastyle}{
	backgroundcolor=\color{backcolour},
	commentstyle=\color{codegray},
	keywordstyle=\color{keywordcolor}\bfseries,
	numberstyle=\tiny\color{codegray},
	stringstyle=\color{codepurple},
	basicstyle=\ttfamily\footnotesize,
	breakatwhitespace=false,
	breaklines=true,
	captionpos=b,
	keepspaces=true,
	numbers=left,
	numbersep=5pt,
	showspaces=false,
	showstringspaces=false,
	showtabs=false,
	tabsize=2,
	language=Java,
	frame=single,
	rulecolor=\color{black!30}
}

\lstset{style=javastyle}

% --- Hinweis-/Beispiel-/Formel-Umgebungen (schlicht, ohne farbige Box) ---
\newenvironment{hinweis}[1][Hinweis]%
	{\par\smallskip\noindent\textbf{#1.}\enspace\ignorespaces}%
	{\par\smallskip}

\newenvironment{beispiel}[1][Beispiel]%
	{\par\smallskip\noindent\textbf{#1:}\par\nobreak\vspace{2pt}}%
	{\par\smallskip}

\newenvironment{formeln}[1][Gegebene Formeln]%
	{\par\smallskip\noindent\textbf{#1:}\par\nobreak\vspace{2pt}}%
	{\par\smallskip}

% ##################################################
% Struktogramm (Nassi-Shneiderman) in reinem TikZ
% ##################################################
\newlength{\nswidth}\setlength{\nswidth}{11.5cm}   % Gesamtbreite
\newlength{\nsind}\setlength{\nsind}{0.55cm}        % Einrueckung je Schleifenebene
% vorberechnete Versatzbreiten fuer Verzweigungen (in Koordinaten nutzbar)
\newlength{\nshalf}\setlength{\nshalf}{\dimexpr\nswidth/2\relax}            % halbe Gesamtbreite
\newlength{\nshalfi}\setlength{\nshalfi}{\dimexpr(\nswidth-\nsind)/2\relax}  % halbe Breite, 1x eingerueckt
\newlength{\nshalfii}\setlength{\nshalfii}{\dimexpr(\nswidth-2\nsind)/2\relax} % halbe Breite, 2x eingerueckt

\tikzset{
	% Basis fuer Anweisungs-/Verzweigungsbloecke
	nb/.style={draw, line width=0.4pt, anchor=north west, inner sep=4pt,
		align=left, font=\footnotesize, fill=white, minimum height=0.8cm},
	% volle Breite
	nbfull/.style={nb, minimum width=\nswidth,
		text width=\dimexpr\nswidth-8pt\relax},
	% eine Ebene eingerueckt (Schleifenrumpf)
	nbi/.style={nb, minimum width=\dimexpr\nswidth-\nsind\relax,
		text width=\dimexpr\nswidth-\nsind-8pt\relax},
	% zwei Ebenen eingerueckt (Schleife in Schleife)
	nbii/.style={nb, minimum width=\dimexpr\nswidth-2\nsind\relax,
		text width=\dimexpr\nswidth-2\nsind-8pt\relax},
	% Verzweigung: halbe Breiten
	nbh/.style={nb, minimum width=\dimexpr\nswidth/2\relax,
		text width=\dimexpr\nswidth/2-8pt\relax},
	nbih/.style={nb, minimum width=\dimexpr(\nswidth-\nsind)/2\relax,
		text width=\dimexpr(\nswidth-\nsind)/2-8pt\relax},
	% Verzweigung halbe Breite, zwei Ebenen eingerueckt
	nbiih/.style={nb, minimum width=\nshalfii,
		text width=\dimexpr\nshalfii-8pt\relax},
	% Kopf einer Verzweigung (das Dreieck), inner sep 0
	nc/.style={draw, line width=0.4pt, anchor=north west, inner sep=0pt,
		minimum height=1.0cm, fill=white},
	ncfull/.style={nc, minimum width=\nswidth},
	nci/.style={nc, minimum width=\dimexpr\nswidth-\nsind\relax},
	ncii/.style={nc, minimum width=\dimexpr\nswidth-2\nsind\relax},
}

% Zeichnet das Dreieck + Beschriftungen in einen bereits gesetzten nc-Knoten.
% #1 Knotenname  #2 Bedingung  #3 ja-Label  #4 nein-Label
\newcommand{\nscond}[4]{%
	\draw[line width=0.4pt] (#1.north west) -- (#1.south);
	\draw[line width=0.4pt] (#1.north east) -- (#1.south);
	\node[font=\scriptsize, anchor=north, inner sep=1pt] at ($(#1.north)+(0,-1.5pt)$) {#2};
	\node[font=\scriptsize, anchor=south west, inner sep=1pt] at ($(#1.south west)+(2.5pt,1.5pt)$) {#3};
	\node[font=\scriptsize, anchor=south east, inner sep=1pt] at ($(#1.south east)+(-2.5pt,1.5pt)$) {#4};
}

% Zeichnet den L-Rahmen einer Schleife (linker + unterer Rand der Einrueckspalte).
% #1 Kopfknoten der Schleife  #2 letzter Rumpfknoten
\newcommand{\nsframe}[2]{%
	\draw[line width=0.4pt] (#1.south west) |- (#2.south west);
}


\title{Probeklausur 4}
\author{Luis Staudt}
\date{}

\begin{document}

% --- Titelblock ---
\begin{center}
	\rule{\textwidth}{0.8pt}\\[0.6em]
	{\LARGE\bfseries Probeklausur 4}\\[0.3em]
	{\large Grundlagen der Programmierung}\\[0.7em]
	\begin{tabular}{r l @{\hspace{1.4cm}} r l}
		\textbf{Studiengänge:}    & PEB \& MVB            & \textbf{Autor:}        & Luis Staudt \\
		\textbf{Semester:}        & Sommersemester 2026   & \textbf{Schwierigkeit:}& mittel \\
		\textbf{Bearbeitungszeit:}& 90 Minuten           & \textbf{Gesamtpunkte:} & 90 \\
	\end{tabular}\\[0.7em]
	\rule{\textwidth}{0.8pt}
\end{center}
\vspace{0.5em}

\noindent\textit{Bearbeiten Sie alle fünf Aufgaben. Achten Sie auf sauberen, kompilierfähigen Java-Quelltext. Hilfsmittel: OpenBook, keine elektronischen Geräte.}

\pagestyle{fancy}
\fancyhf{}
\lhead{\small Probeklausur 4, Grundlagen der Programmierung}
\rhead{\small SoSe 2026 (mittel)}
\cfoot{\thepage}
\renewcommand{\headrulewidth}{0.4pt}

% =============================================
\clearpage
\section*{Aufgabe 1: Struktogramm in Java umsetzen \hfill (24 Punkte)}

Das folgende Struktogramm prüft zunächst eine Eingabe, bildet danach die Collatz-Folge eines Startwerts (und merkt sich deren größten erreichten Wert) und berechnet zum Schluss eine Fakultät mit Überlaufkontrolle.

Setzen Sie das Struktogramm \textbf{vollständig in ein lauffähiges Java-Programm} um.

\begin{center}
\begin{tikzpicture}
	\node[nbfull] (a1) at (0,0) {lies Startwert $k$ ein};

	% Eingabepruefung
	\node[nbfull] (h0) at (a1.south west) {\textbf{solange} $k < 1$};
	\node[nbi]    (v0) at ($(h0.south west)+(\nsind,0)$) {lies $k$ erneut ein};
	\nsframe{h0}{v0}

	\node[nbfull] (a2) at ($(v0.south west)+(-\nsind,0)$) {\texttt{schritte = 0;}\qquad \texttt{maxK = k;}};

	% Collatz
	\node[nbfull] (h1) at (a2.south west) {\textbf{solange} $k \neq 1$};
	\node[nci]    (c1) at ($(h1.south west)+(\nsind,0)$) {};
	\nscond{c1}{\texttt{k \% 2 == 0}}{ja}{nein}
	\node[nbih]   (c1l) at (c1.south west) {\texttt{k = k / 2;}};
	\node[nbih]   (c1r) at ($(c1.south west)+(\nshalfi,0)$) {\texttt{k = 3 * k + 1;}};
	\node[nci]    (c2) at (c1l.south west) {};
	\nscond{c2}{\texttt{k > maxK}}{ja}{nein}
	\node[nbih]   (c2l) at (c2.south west) {\texttt{maxK = k;}};
	\node[nbih]   (c2r) at ($(c2.south west)+(\nshalfi,0)$) {\textit{(nichts)}};
	\node[nbi]    (b1) at (c2l.south west) {\texttt{schritte = schritte + 1;}};
	\nsframe{h1}{b1}

	\node[nbfull] (a3) at ($(b1.south west)+(-\nsind,0)$) {drucke \texttt{schritte} und \texttt{maxK}};
	\node[nbfull] (a4) at (a3.south west) {lies $n$ ein;\qquad \texttt{fak = 1;}};

	% Fakultaet
	\node[nbfull] (h2) at (a4.south west) {\textbf{für} $i = 2,\,3,\,\ldots,\,n$};
	\node[nbi]    (f1) at ($(h2.south west)+(\nsind,0)$) {\texttt{fak = fak * i;}};
	\node[nci]    (c3) at (f1.south west) {};
	\nscond{c3}{\texttt{fak > grenze}}{ja}{nein}
	\node[nbih]   (c3l) at (c3.south west) {drucke $i$, Warnung; Schleife abbrechen};
	\node[nbih]   (c3r) at ($(c3.south west)+(\nshalfi,0)$) {\textit{(nichts)}};
	\nsframe{h2}{c3l}

	\node[nbfull] (a5) at ($(c3l.south west)+(-\nsind,0)$) {drucke \texttt{fak}};
\end{tikzpicture}
\end{center}

% =============================================
\clearpage
\section*{Aufgabe 2: Kinematik, gesuchte Größe auswählen \hfill (14 Punkte)}

Bei der gleichmäßig beschleunigten Bewegung kann je nach Aufgabenstellung eine andere Größe gesucht sein.

\begin{formeln}
\[
	\text{(1)}\;\; v = v_0 + a\,t
	\qquad
	\text{(2)}\;\; s = v_0\,t + \tfrac{1}{2}\,a\,t^2
\]
\[
	\text{(3)}\;\; a = \frac{v - v_0}{t}
	\qquad
	\text{(4)}\;\; t = \frac{v - v_0}{a}
\]
\end{formeln}

Schreiben Sie ein Programm, das die gesuchte Größe per Zahl (\texttt{1} bis \texttt{4}) wählen lässt, die jeweils benötigten Werte abfragt und das Ergebnis mit der \textbf{richtigen} Formel berechnet.
Fangen Sie die Division durch Null ab: bei Auswahl~\texttt{3} den Fall $t = 0$, bei Auswahl~\texttt{4} den Fall $a = 0$.

\begin{beispiel}[Beispieldialog]
\begin{lstlisting}[language={}]
Gesuchte Groesse (1=v, 2=s, 3=a, 4=t): 1
v0 [m/s]: 5
a [m/s^2]: 2
t [s]: 3
Geschwindigkeit v = 11.0 m/s
\end{lstlisting}
\end{beispiel}

% =============================================
\clearpage
\section*{Aufgabe 3: Punkt drehen und verschieben \hfill (20 Punkte)}

Gegeben ist die Klasse \texttt{Punkt} (Attribute \texttt{xKoordinate}, \texttt{yKoordinate} als \texttt{double}).
Implementieren Sie zwei Methoden, die jeweils einen \textbf{neuen} \texttt{Punkt} zurückgeben (die aufrufende Instanz bleibt unverändert):

\begin{itemize}[nosep]
	\item \texttt{rotiere(double winkel)}, dreht den Punkt um den Ursprung (Winkel in Grad).
	\item \texttt{verschiebeUm(double strecke, double richtung)}, verschiebt den Punkt um \texttt{strecke} in Richtung \texttt{richtung} (in Grad).
\end{itemize}

\begin{formeln}[Berechnungsformeln (Winkel in Grad, daher Umrechnung $\cdot\,\pi/180$)]
\[
	x_n = x\cos\beta - y\sin\beta \qquad y_n = x\sin\beta + y\cos\beta
\]
\[
	x_n = x + \text{strecke}\cdot\cos(\text{richtung}) \qquad
	y_n = y + \text{strecke}\cdot\sin(\text{richtung})
\]
\end{formeln}

\begin{lstlisting}
Punkt p = new Punkt();
p.xKoordinate = 3;  p.yKoordinate = 4;

Punkt gedreht  = p.rotiere(90);          // um 90 Grad drehen
Punkt versetzt = p.verschiebeUm(5, 0);   // 5 Einheiten in x-Richtung
System.out.println(gedreht.xKoordinate + " / " + gedreht.yKoordinate);
\end{lstlisting}

% =============================================
\clearpage
\section*{Aufgabe 4: Fachwerk, längster und kürzester Stab \hfill (16 Punkte)}

Gegeben ist die aus Vorlesung und Übung bekannte Datei \texttt{fachwerk.jar}.
Für ein gewähltes Fachwerk sollen der \textbf{längste} und der \textbf{kürzeste} Stab (jeweils mit Stabnummer und Länge) sowie die \textbf{Gesamtlänge} aller Stäbe ausgegeben werden.

\begin{formeln}[Stablänge aus den Knotenkoordinaten]
\[ l = \sqrt{(x_e - x_a)^2 + (y_e - y_a)^2} \]
\end{formeln}

\begin{center}
\begin{tikzpicture}[>=Stealth, scale=1.0,
	knoten/.style={circle, fill=black, inner sep=1.6pt},
	stab/.style={thick, gray!70!black}]
	\coordinate (A) at (0,0); \coordinate (B) at (2.5,0); \coordinate (C) at (5,0);
	\coordinate (D) at (1.0,2.0); \coordinate (E) at (4.0,2.4);
	\draw[stab] (A)--(B); \draw[stab] (B)--(C);
	\draw[stab] (A)--(D); \draw[stab] (D)--(B); \draw[stab] (D)--(E);
	\draw[stab] (E)--(B); \draw[stab] (E)--(C);
	\foreach \p in {A,B,C,D,E} \node[knoten] at (\p) {};
\end{tikzpicture}\\[-0.3em]
{\footnotesize\itshape Schematische Darstellung eines Fachwerks.}
\end{center}

\begin{hinweis}[Wichtig]
	Es geht nur um die \emph{Geometrie} des Fachwerks (die Knotenkoordinaten).
	Ein Aufruf von \texttt{berechneGroessen()} ist daher \textbf{nicht} erforderlich.
	Benötigt: \texttt{gibAnzahlStaebe()}, \texttt{gibStab(i)}, \texttt{gibKnotenNrA()}, \texttt{gibKnotenNrE()}, \texttt{gibKnoten(i)}, \texttt{gibX()}, \texttt{gibY()}.
\end{hinweis}

\begin{lstlisting}
public class Fachwerkaufgabe {
    public static void main(String[] args) {
        System.out.println("Welches Fachwerk wollen Sie rechnen (0-4)?");
        int wahl = Tastatureingabe.leseInt();
        Fachwerk fw = new Fachwerk(wahl);
        // hier erweitern: laengsten/kuerzesten Stab und Gesamtlaenge bestimmen
    }
}
\end{lstlisting}

\begin{beispiel}[Mögliche Ausgabe]
\begin{lstlisting}[language={}]
Laengster Stab:  Stab 7 mit Laenge 4.272
Kuerzester Stab: Stab 2 mit Laenge 2.0
Gesamtlaenge aller Staebe: 35.61
\end{lstlisting}
\end{beispiel}

% =============================================
\clearpage
\section*{Aufgabe 5: Energie-Wertetabelle in eine Datei schreiben \hfill (16 Punkte)}

Ein Körper der Masse $m = 3~\text{kg}$ wird aus der Ruhe mit $a = 2~\text{m/s}^2$ beschleunigt.
Für Geschwindigkeit und kinetische Energie gilt
\[ v(t) = a\,t \qquad E(t) = \tfrac{1}{2}\,m\,\bigl(a\,t\bigr)^2. \]

Schreiben Sie ein Programm, das für $t = 0~\text{s}$ bis $t = 10~\text{s}$ in Schritten von $0{,}5~\text{s}$ jeweils $t$, $v$ und $E$ als drei Spalten in die Datei \texttt{energie.txt} schreibt. Hängen Sie am Ende eine Zeile an, die den \textbf{Zeitpunkt} nennt, zu dem $E$ erstmals eine Grenze $E_{\text{grenze}} = 100~\text{J}$ überschreitet.

\begin{hinweis}[FileWriter]
	\lstinline|new FileWriter("energie.txt")|, je Zeile \texttt{write(...)} mit \lstinline|"\n"|, am Ende \texttt{close()}.
	Nutzen Sie \lstinline|"\t"| für die Spaltenabstände.
\end{hinweis}

\begin{beispiel}[Erwarteter Dateiinhalt (Auszug)]
\begin{lstlisting}[language={}]
t [s]   v [m/s]   E [J]
0.0     0.0       0.0
0.5     1.0       1.5
...
6.0     12.0      216.0
...
Grenze von 100 J erstmals ueberschritten bei t = 4.5 s
\end{lstlisting}
\end{beispiel}

\end{document}
